\subsubsection{Access methods}
\inlinedefinition[Access Methods] typically refers to the operations and data structures used for reading data from the buffer cache into the private space of the worker thread

The basic access method is \bi{sequential scan}, where all pages of the table are read from beginning to end.
Indexing is used to avoid having to scan the whole table.

There are many indexing techniques, but we will focus on three indexing techniques in this section, hashing, B+ trees and bitmaps.

Hashing is primarily used to randomize access to the data, finding data based on a key (hash and access concept)
and both as an index for tables and used to implement operators.

B+ trees on the other hand are used to order data along the values of one attribute, but require traversal, unless we are looking for consecutive values.
They are thus only used as an indexing mechanism.
